\section{Related Work}
\label{sec:relatedwork}

\paragraph{Large Language Models}
LLMs have shown impressive performances across various tasks~\cite{GPT-4, Gemini}, including scientific fields such as mathematics, physics, medicine, and computer science~\cite{author/discovery, MathematicalLLM, LLM/discovery/2, Airesearchagent, idea/co-creation}. For instance, GPT-4 can understand DNA sequences, design biomolecules, predict molecular behavior, and solve PDE problems~\cite{LLM/discovery/1}. However, LLMs have mainly been used for accelerating the experimental validation of already identified research ideas, but not for identifying new problems. 

\paragraph{Hypothesis Generation}
The principle of hypothesis generation is based on literature-based discovery~\cite{LBD}, which aims to discover relationships between concepts~\cite{Henry2017LiteratureBD}. For instance, these concepts could be a specific disease and a compound not yet considered as a treatment for it. Early works on automatic hypothesis generation first build a corpus of discrete concepts, and then identify their relationships with machine learning approaches, e.g., using similarities between word (concept) vectors~\cite{hypothesis/wordvector} or applying link prediction methods over a graph (where concepts are nodes)~\cite{hypothesis/graph/1, hypothesis/graph/2}. Recent approaches are further powered by LLMs~\cite{CLBD, zeroshot/hypothesis, tomato}, leveraging their prior knowledge about scientific disciplines. Yet, all these approaches perform idea generation in a localized manner and are designed to identify potential relationships between two variables or generate sentence-level connections, which may be sub-optimal to capture the complexity and multifaceted nature of real-world problems (e.g., urban planning involves numerous interacting variables). Meanwhile, we do not artificially restrict the target research idea to be a predictive single concept or simple binary link, instead allowing the model to generate ideas in a more open-ended fashion.

We note that there has been a recent surge of interest in exploring scientific idea generation: from \citet{li2024chainideasrevolutionizingresearch} that focus on evaluating whether LLMs can generate research ideas that are better than human ideas, to \citet{lu2024aiscientistfullyautomated} that aim to automatically generate full research papers (including idea development, code writing, and experiment execution), to \citet{li2024chainideasrevolutionizingresearch} that enhance the idea generation process by organizing a sequential chain of literature, all of which acknowledge and build upon insights from a prior version of our paper.

\paragraph{Knowledge-Augmented LLMs}
The approach to augment LLMs with external knowledge makes them more accurate and relevant to target contexts. Much prior work aims at improving the factuality of LLM responses to queries by retrieving the relevant documents and injecting them into the LLM input~\cite{InternetaugmentedLM, InContextRetrieval, REPLUG}. In addition, given that entities or facts are atomic units for representing knowledge, recent studies augment LLMs with them~\cite{KAPING, RetrieveRewriteAnswer}. In contrast to these efforts, which use knowledge units piecemeal, we instead jointly leverage accumulated knowledge over massive troves of scientific papers. Also, \citet{K-LaMP} proposes to use entities for query suggestion, which -- while similar -- has the entirely different objective of narrowing the focus of LLMs to entities already present in their context. Instead, our approach retrieves and integrates entities outside the given context, enabling LLMs to explore other concepts for research idea generation.

\paragraph{Iterative Refinements with LLMs}
Similar to humans, LLMs do not always generate optimal outputs on their first attempt. To tackle this, drawing inspiration from humans who can iteratively refine their thoughts based on critiques from themselves and their peers, many recent studies have investigated the potential of LLMs to correct and refine their outputs, demonstrating that they indeed possess those capabilities~\cite{self-refine/0, self-refine/1, self-refine/2, self-refine/3, CLBD, zeroshot/hypothesis, tomato}. Based on their findings, we extend this paradigm (and further test their capability) to our novel scenario of research idea generation.


